\section{Kerr geometry and spatial microlocal separation}
\label{sec:geometry}

The exterior construction requires a spacelike stationary slicing, a
description of all null rays that can remain in a bounded cylinder, and a
spatial conic neighborhood of trapping on which the stationary operator is
elliptic at zero frequency.  We establish these properties for every fixed
subextremal Kerr metric.  Although the mass has been normalized to $M=1$ in
the main argument, we retain $M$ in the identities below to display their
parameter dependence.  Throughout this section,
\[
 M>0,\qquad 0<\mathfrak a<M
\]
are fixed.  Constants may depend on this pair; no assertion is uniform as
$\mathfrak a\to0$ or $\mathfrak a\to M$.

\subsection{The stationary slicing and its ends}

At either exterior end, let $T$ be the future Boyer--Lindquist time.  In the
trapped block we continue the corresponding sign of the Killing time,
reversing the axial coordinate when necessary.  Write
\begin{equation}
 \begin{gathered}
 \Sigma_K=r^2+\mathfrak a^2\cos^2\theta,\qquad
 \Delta=r^2-2Mr+\mathfrak a^2=(r-r_+)(r-r_-),\\
 r_\pm=M\pm\sqrt{M^2-\mathfrak a^2},\qquad
 r_*'=\frac{r^2+\mathfrak a^2}{\Delta},\qquad
 \kappa_\pm=\frac{r_+-r_-}{2(r_\pm^2+\mathfrak a^2)}.
 \end{gathered}
 \label{geo:kerr-functions}
\end{equation}
In particular, $r_->0$ and both numbers $\kappa_\pm$ are positive; at the
inner horizon we use the absolute value of the surface gravity.

The exterior computing cylinder will extend to an inner face $r=r_b$ with
$r_-<r_b<r_+$.  Its precise depth is chosen after the entry cylinder for the
double-null argument and its finite margins have been fixed.  The inner
face is placed strictly deeper than that entry cylinder.

\begin{lemma}[A bent spacelike time]
\label{geo:slicing}
For any prescribed finite interval $[r_b,R_1]$, where $R_1>r_+$,
one can choose a stationary time
\begin{equation}
 t=T+r_*-G(r)
 \label{geo:time}
\end{equation}
with the following properties.  One has $G'=1$ on the prescribed region,
$1\leq G'\leq r_*'$ in a subsequent exterior interpolation, and $G=r_*$
near infinity.  The function $t$ is regular at the future event horizon,
its level hypersurfaces are spacelike, and $r=r_b$ is strictly future
outflow.  The delay $G-r_*$ at each fixed depth above $r_-$ can be made
arbitrarily large by moving the exterior bend outward.
\end{lemma}

\begin{proof}
Use the regular ingoing azimuth near the horizon,
\[
 d\widetilde\varphi=d\varphi+\frac{\mathfrak a}{\Delta}\,dr,
\]
and return smoothly to the Boyer--Lindquist azimuth toward infinity.  These
coordinates, together with $r$ and local sphere coordinates, give
stationary spatial coordinates $y$; near infinity we use their
Boyer--Lindquist Cartesian version.  Let $b$ denote the Kerr metric in
$(t,y)$.  Direct substitution in the Boyer--Lindquist inverse metric gives
\begin{equation}
 \Sigma_K b^{-1}(dt,dt)
 =-2(r^2+\mathfrak a^2)G'+\Delta(G')^2
      +\mathfrak a^2\sin^2\theta.
 \label{geo:time-norm}
\end{equation}
For $G'=1$, the right-hand side is
\begin{equation}
 -\Sigma_K-2Mr<0.
 \label{geo:time-norm-one}
\end{equation}
In the exterior it is a convex quadratic in $G'$.  Its value at
$G'=r_*'=(r^2+\mathfrak a^2)/\Delta$ is
\[
 -\frac{(r^2+\mathfrak a^2)^2}{\Delta}
       +\mathfrak a^2\sin^2\theta<0.
\]
Both endpoint values are negative, so every indicated interpolation is
spacelike.  Near the horizon, $t$ is the regular ingoing time minus the
smooth function $G$.  Between the horizons $b^{-1}(dr,dr)=\Delta/\Sigma_K<0$;
with the chosen future orientation, $r$ decreases along future causal
curves.  This proves strict outflow at $r_b$.

Choose the additive constant of $G$ by imposing $G=r_*$ after the
interpolation.  Extending the interval where $G'=1$ changes the delay at
any fixed smaller radius by the integral of $r_*'-1$ over the added
interval.  Since
\[
 r_*'-1=\frac{2Mr}{\Delta}=\frac{2M}{r}+O(r^{-2}),
\]
this integral tends to positive infinity as the bend moves outward.  Thus
the delay can satisfy all fixed-depth margins at once, while $G'=1$
continues to hold throughout the prescribed interval.
\end{proof}

As in \Cref{cmp:bent-time}, the delays place the $t=0$ computing slices on
or to the future of the bridge and separate their trapped-block portions
by the required finite margins.  The order of choices is important:
first choose the entry depth and these margins, then choose $r_b$, and
finally move the bend far enough outward.  Lemma~\ref{geo:slicing} permits
this order without altering the slicing on the prescribed interval.
After proving compactness of the trapping radii below, we take $R_1$
larger than all of them.

Near infinity $t=T$, so all components and inverse components have the
ordinary stationary Kerr symbol expansions.  In particular,
\begin{equation}
 \begin{gathered}
 b-\eta=O(r^{-1}),\qquad b^{ti}=O(r^{-2}),\qquad
 \Gamma(b)=O(r^{-2}),\\
 r_*=r+2M\log r+O(1),\\
 b_{tt}=-1+\frac{2M}{r}+O(r^{-2}),\qquad
 b_{ij}=\delta_{ij}+\frac{2M y_i y_j}{r^3}+O(r^{-2}).
 \end{gathered}
 \label{geo:end-expansions}
\end{equation}
Here $\eta$ is the Minkowski metric, and every remainder includes its
differentiated symbol bounds.  The terms of order $r^{-1}$ are rotation
invariant.  The leading term of $b_{ti}$ is a nonzero constant multiple of
$(\mathbf J\times y)_i/r^3$, where
$\mathbf J=M\boldsymbol{\mathfrak a}$ and
$\boldsymbol{\mathfrak a}$ is the oriented spin vector in the chosen
axes.  These assertions follow by expanding the Boyer--Lindquist metric;
the angular expressions and their rotations are smooth on the sphere.

\subsection{Radial dynamics and the horizon source}

We now classify the null rays in a large computing cylinder.  A future
null covector means a nonzero null covector $k$ whose metric dual
$k^\sharp$ is future directed.  Set
\begin{equation}
 \begin{gathered}
 E=-k(\partial_t)=-k(\partial_T),\qquad
 L_z=k(\partial_\varphi),\\
 D=k_\theta^2+
     \left(\frac{L_z}{\sin\theta}
                    -\mathfrak a E\sin\theta\right)^2,\qquad
 P=(r^2+\mathfrak a^2)E-\mathfrak a L_z.
 \end{gathered}
 \label{geo:constants}
\end{equation}
The angular expression is understood intrinsically at the poles: its
expansion is the smooth round cotangent norm, minus
$2\mathfrak a E L_z$, plus $\mathfrak a^2E^2\sin^2\theta$.
The separated principal symbol in Boyer--Lindquist variables is
\begin{equation}
 q_E:=\Sigma_K b^{-1}(k,k)
       =\Delta k_r^2+D-\frac{P^2}{\Delta}.
 \label{geo:separated-symbol}
\end{equation}
Stationarity, axial symmetry, and separation show that $E,L_z,D$ are
constant along null rays.  With a dot denoting the affine geodesic
parameter, the corresponding equations are
\begin{equation}
 \begin{aligned}
 (\Sigma_K\dot r)^2&=P^2-\Delta D,\\
 \Sigma_K\dot T&=
  \frac{(r^2+\mathfrak a^2)P}{\Delta}
       +\mathfrak a(L_z-\mathfrak a E\sin^2\theta).
 \end{aligned}
 \label{geo:ray-equations}
\end{equation}
These are Carter's separated geodesic equations
\cite{Carter1968}, in the conventions used for the characteristic-flow
calculation below.

\begin{lemma}[Radial alternatives]
\label{geo:radial}
For the fixed parameters above, future null rays in $r>r_+$ have $P>0$.
If $E\leq0$, their allowed radial set has at most one simple turning
point, and there is no photon trapping.  If $E>0$, every trapped ray is a
spherical photon ray, with constant radius satisfying
\begin{equation}
 F=\frac{2r\Delta}{r-M},\qquad
 \frac{D}{E^2}=\frac{4r^2\Delta}{(r-M)^2},\qquad
 F:=\frac PE=r^2+\mathfrak a^2-\mathfrak a\frac{L_z}{E}.
 \label{geo:double-root}
\end{equation}
The energy-normalized trapped directions form a compact set separated
from $r=r_+$ and from infinity.  The relevant radial critical point is a
unique nondegenerate minimum.  A simple radial turning point reverses
toward the same end of its allowed interval, and there is no homoclinic
orbit from a double radial root back to that root.
\end{lemma}

\begin{proof}
In the exterior,
\[
 X=(r^2+\mathfrak a^2)\partial_T+\mathfrak a\partial_\varphi
\]
is future timelike: direct evaluation gives
$b(X,X)=-\Sigma_K\Delta$.  Since $k(X)=-P$, pairing with the future
null vector $k^\sharp$ gives $P>0$.

Suppose first that $E\leq0$.  Then $P'=2rE\leq0$, while
$\sqrt{\Delta D}$ is strictly increasing when $D>0$.  The case $D=0$
cannot occur: it would force $k_\theta=0$ and
$L_z=\mathfrak a E\sin^2\theta$, and hence $P=E\Sigma_K\leq0$.
Thus the inequality $P^2\geq\Delta D$, on its positive-$P$ branch,
has at most one boundary point, and that point is simple.  It cannot
contain either a double root or a bounded allowed island.

For $E>0$, restrict to the interval where $F>0$.  Differentiation gives
\begin{equation}
 \left(\frac{F}{\sqrt\Delta}\right)'
       =\frac{N}{\Delta^{3/2}},\qquad
 N=2r\Delta-(r-M)F,
 \label{geo:radial-derivative}
\end{equation}
and
\begin{equation}
 N'=2\Delta+2r(r-M)-F,\qquad N''=6(r-M)>0.
 \label{geo:radial-convexity}
\end{equation}
If $F(r_+)>0$, then $N(r_+)<0$.  The strictly convex function $N$
tends to positive infinity and has exactly one zero beyond $r_+$;
its derivative at that zero is strictly positive.  Consequently
$F/\sqrt\Delta$ has a unique nondegenerate minimum.  If instead
$F(r_+)\leq0$, let $r_0\geq r_+$ be its first zero.  Since
$F'=2r>0$,
\[
 N(r_0)=2r_0\Delta(r_0)\geq0,\qquad
 N'(r_0)=2\Delta(r_0)+2r_0(r_0-M)>0.
\]
Equation~\eqref{geo:radial-convexity} then gives $N>0$ wherever
$F>0$.  In this second case there is no radial critical point.

A double root of $P^2-\Delta D$ must therefore occur at the unique
minimum in the first case.  Setting $N=0$ and
$F^2=\Delta D/E^2$ gives \eqref{geo:double-root}.  Moreover,
\begin{equation}
 \left|\frac{L_z}{E}\right|-\mathfrak a
                  \leq\frac{\sqrt D}{E}.
 \label{geo:impact-bound}
\end{equation}
If double roots approached $r_+$, their identities would force
\[
 \frac{L_z}{E}\longrightarrow
       \frac{r_+^2+\mathfrak a^2}{\mathfrak a}>\mathfrak a,
 \qquad \frac{D}{E^2}\longrightarrow0,
\]
contrary to \eqref{geo:impact-bound}.  They cannot approach infinity
either: \eqref{geo:double-root} makes $|L_z/E|$ quadratic in $r$ to
leading order, whereas $\sqrt D/E$ is only linear.  The radial range,
the two impact quantities, and the normalized angular covectors are
therefore bounded.  Smoothness of the angular expression at the poles
gives compactness in the intrinsic cotangent bundle.

Finally, the same one-minimum description classifies the allowed radial
intervals.  A simple endpoint is a reflection toward the same end of its
interval.  At a double root, $F^2/\Delta$ increases strictly in either
radial direction away from the minimum, so a departing branch cannot
turn back to it.  There are no further radially bounded trajectories.
Covector reversal gives the corresponding description of the past
sheet.
\end{proof}

\begin{lemma}[The outer-horizon radial source]
\label{geo:horizon-source}
The outer-horizon generators are a strict radial source for the null
ray flow parameterized by increasing regular time $t$ on spatial
cosphere directions.  More precisely, let $H$ be the horizon Killing
generator, normalized by its Killing parameter, so that
\[
 \nabla_HH=\kappa_+H,\qquad H^t>0.
\]
There is a smooth positive degree-one frequency scale $b_1$, equal to
$|\rho|$ on $k=\rho H^\flat$, such that, for the increasing-time ray
field $V$,
\begin{equation}
 V\log b_1=-\frac{\kappa_+}{H^t}
       \quad\hbox{on the generators}.
 \label{geo:source-frequency}
\end{equation}
On null directions near the generators, the degree-zero function
\begin{equation}
 f_H=\frac{E^2+D}{b_1^2}
 \label{geo:source-function}
\end{equation}
vanishes precisely at the generators and satisfies
$Vf_H\geq c_H f_H$ for some $c_H>0$ after shrinking the neighborhood.
\end{lemma}

\begin{proof}
The only null tangencies to a null horizon are its generators.
Along an affine generator, write $k^\sharp=\rho H$ and use a
Killing-invariant normalization of $H$.  The affine geodesic equation
and $\nabla_HH=\kappa_+H$ give
\[
 H(\log|\rho|)=-\kappa_+.
\]
Dividing by $H^t$ proves \eqref{geo:source-frequency}.  Extend
$|\rho|$ smoothly and homogeneously to a positive frequency scale in a
conic neighborhood.  Since $E$ and $D$ are conserved,
\[
 Vf_H=-2(V\log b_1)f_H.
\]
The coefficient is strictly positive on the generators and remains
bounded below nearby.  On the null set, $E=D=0$ forces
$L_z=k_\theta=0$; near the horizon the remaining null direction is
exactly its generator covector.  This proves the asserted zero set
and the strict source inequality.
\end{proof}

Together, Lemmas~\ref{geo:radial} and \ref{geo:horizon-source} give the
escape alternatives used below.  Between the horizons $dr$ is timelike,
so every future ray there moves toward the inner outflow face.  In a
sufficiently large exterior cylinder, a ray away from trapping and the
generators escapes in at least one time direction, or approaches the
generator source in backward regular time.  An orbit approaching
$r=r_+$ without a transverse crossing has a normalized invariant limit
set tangent to that horizon, and hence consisting of generators.
The radial alternatives exclude any additional recurrent set.

\subsection{Normal hyperbolicity and tangent dynamics}

The radial classification identifies the only exterior obstruction to
escape.  The smooth symplectic structure and normal hyperbolicity of
the subextremal Kerr trapped set are established in
\cite[Section~II]{DyatlovZworski2013}.  For the estimates below we give
explicit defining functions for its stable and unstable surfaces and
verify that their exponential rates dominate every fixed-order
tangent derivative of the trapped flow.

\begin{lemma}[Transverse defining functions]
\label{geo:normal-functions}
Near the energy-normalized trapped set on either temporal sheet, the
critical radius $r_c=r_c(L_z/E)$ is smooth.  There are smooth homogeneous
degree-zero functions
\begin{equation}
 \begin{aligned}
 z_1&=\frac{\Sigma_K\dot r}{E\sqrt\Delta},\\
 z_2&=\operatorname{sgn}(r-r_c)
       \sqrt{\frac{F(r)^2}{\Delta(r)}
                  -\frac{F(r_c)^2}{\Delta(r_c)}}
 \end{aligned}
 \label{geo:normal-variables}
\end{equation}
whose common zero set is trapping.  For the increasing-time ray field,
\begin{equation}
 Vz_2=\lambda z_1,\qquad Vz_1=\lambda z_2,
       \qquad \lambda>0.
 \label{geo:normal-system}
\end{equation}
Consequently $\varphi_\pm=z_1\pm z_2$ satisfy
$V\varphi_\pm=\pm\lambda\varphi_\pm$.  Their reduced spatial
symplectic bracket on either halfwave is nonzero on the trapped set.
The transverse rates have a uniform positive lower bound on each
fixed normalized trapped set.
\end{lemma}

\begin{proof}
The strict minimum established in Lemma~\ref{geo:radial} and the
implicit function theorem give the smooth function $r_c$.  Set
$U=F^2/\Delta$ and $U_c=U(r_c)$.  Near the minimum,
\[
 U(r)-U_c=(r-r_c)^2 a_0(r,L_z/E),\qquad a_0>0
\]
with $a_0$ smooth.  Thus $z_2=(r-r_c)\sqrt{a_0}$ is smooth and
$\partial_rz_2>0$ after restricting the neighborhood.  On the null
set, \eqref{geo:ray-equations} gives
\begin{equation}
 z_1^2-z_2^2=U_c-\frac{D}{E^2},
 \label{geo:normal-invariant}
\end{equation}
which is conserved.  Since $E$ and $L_z$ are constant,
\[
 Vz_2=(\partial_rz_2)\frac{\dot r}{\dot t}=\lambda z_1,
 \qquad
 \lambda:=(\partial_rz_2)\frac{E\sqrt\Delta}{\Sigma_K\dot t}.
\]
On a trapped future ray $E>0$, while $\dot t>0$ because the slicing
is spacelike.  Both signs reverse on the opposite sheet, so
$E/\dot t>0$ on both sheets.  Hence $\lambda>0$ there and nearby.
Differentiating \eqref{geo:normal-invariant} gives the second equation
of \eqref{geo:normal-system} where $z_1\neq0$; smoothness extends it
across $z_1=0$.  Compactness gives the positive lower bound.

For symplectic transversality first hold $E$ fixed in the
Boyer--Lindquist spatial phase space.  Then
$z_1=\sqrt\Delta\,k_r/E$, and the radial contribution to
$\{z_1,z_2\}$ is, up to the bracket convention,
\[
 \frac{\sqrt\Delta}{E}\,\partial_rz_2\neq0.
\]
There is no axial contribution because the functions are independent
of the axial coordinate, and no polar contribution.  At fixed energy,
a stationary change of time adds an exact one-form times $E$ to the
spatial covector, while a spatial coordinate change acts by cotangent
pullback.  The reduced bracket is therefore unchanged.

Finally pass from fixed $E$ to a temporal root $E=e(y,\eta)$ of the
characteristic equation.  Write $f|_e=f(y,\eta,e(y,\eta))$, and use
the bracket convention $H_f g=\{f,g\}$.  Implicit differentiation
gives, on the characteristic set,
\begin{equation}
 \{f|_e,g|_e\}
 =\{f,g\}
   +\frac{g_E\{q_E,f\}-f_E\{q_E,g\}}{\partial_E q_E},
 \qquad \partial_Eq_E=-2\Sigma_K k^\sharp(dt)\neq0,
 \label{geo:root-bracket}
\end{equation}
where the symbol is expressed in the regular stationary coordinates,
the brackets on the right hold $E$ fixed, and $\partial_E$ holds
their spatial covector $\eta$ fixed.  The denominator is nonzero
because $dt$ is timelike.  For $f=z_1$ and $g=z_2$, both
characteristic derivatives in the numerator vanish on the invariant
trapped set by \eqref{geo:normal-system}.  Thus the nonzero
fixed-energy bracket survives on each halfwave.  Since
$\{\varphi_+,\varphi_-\}=\pm2\{z_1,z_2\}$, the assertion follows.
\end{proof}

\begin{lemma}[Polynomial growth of tangent derivatives]
\label{geo:tangent-growth}
Fix $E\neq0$.  On a neighborhood of the null trapped set in the radial
critical manifold
\[
 r=r_c(L_z/E),\qquad k_r=0,
\]
the flow of the separated fixed-frequency symbol has polynomial
growth of every fixed-order tangent derivative.  The bounds are
uniform on compact invariant subsets, include variation through a
sufficiently small interval of nearby symbol levels, and persist
under positive smooth time changes with speed bounded above and
away from zero.  In particular, they hold in the regular stationary
time $t$, on either sheet with its increasing-time orientation.
\end{lemma}

\begin{proof}
We first establish a strict angular margin on the null trapped set.
There $D=P^2/\Delta$ and $\Delta<r^2$.  Expanding $P$ gives
\begin{equation}
 \begin{aligned}
 D+2\mathfrak a E L_z
 &>\frac{P^2}{r^2}+2\mathfrak a E L_z\\
 &=(r^2+2\mathfrak a^2)E^2
       +\frac{\mathfrak a^2(\mathfrak a E-L_z)^2}{r^2}
  \geq(r^2+2\mathfrak a^2)E^2.
 \end{aligned}
 \label{geo:angular-margin}
\end{equation}
Thus
\begin{equation}
 c_1=D+2\mathfrak a E L_z-2\mathfrak a^2E^2>0.
 \label{geo:oscillator-coefficient}
\end{equation}
Compactness makes this inequality uniform, and it persists on nearby
levels.  Also
\[
 |p_{\mathbb S^2}|^2
  =D+2\mathfrak a E L_z-\mathfrak a^2E^2\sin^2\theta
  >\mathfrak a^2E^2,
\]
so the angular momentum covector never vanishes.

Use half the Hamiltonian \eqref{geo:separated-symbol} to normalize
the separated time.  On the radial critical manifold its angular
field is $H_{D/2}+\omega(L_z)H_{L_z}$, where $\omega$ is smooth
on the relevant impact interval.  The second term is an axial
rotation and preserves $D,L_z$.  Put
\[
 A=\mathfrak a^2E^2,\qquad
 C=D+2\mathfrak a E L_z,\qquad
 Q=C-L_z^2-A=D-(L_z-\mathfrak a E)^2,
\]
and let $x_1=\cos\theta$.  The globally smooth function $x_1$
satisfies
\begin{equation}
 \begin{aligned}
 (x_1')^2&=(1-x_1^2)C-L_z^2-A(1-x_1^2)^2,\\
 (x_1')^2+c_1x_1^2+Ax_1^4&=Q,\qquad
 x_1''=-c_1x_1-2Ax_1^3.
 \end{aligned}
 \label{geo:oscillator}
\end{equation}
The second-order equation follows by differentiation away from the
turning points and extends to them by smoothness.  It is a confining
oscillator because $c_1>0$.

In particular $Q\geq0$, and its squared amplitude is
\begin{equation}
 d_1^2=\frac{2Q}{c_1+\sqrt{c_1^2+4AQ}}.
 \label{geo:amplitude}
\end{equation}
This formula is smooth even at $Q=0$.  Substituting
$x_1=d_1\sin u$ gives the positive invariant period
\begin{equation}
 T_{\mathrm{ang}}(D,L_z)
   =\int_0^{2\pi}
       \frac{du}{\sqrt{c_1+Ad_1^2(1+\sin^2u)}}.
 \label{geo:angular-period}
\end{equation}
It depends smoothly on the conserved quantities, including at the
constant equatorial orbit where its value is $2\pi/\sqrt{c_1}$.
It also remains smooth when the orbit passes through a pole.

Let $\Phi^s$ denote this angular flow, and let $R_\zeta$ be axial
rotation by $\zeta\in\mathbb S^1$.  The circle action is free on
the angular phase states under consideration.  Away from the poles
it is free on the base; at a pole a nonzero covector has trivial
rotational stabilizer.  The data $D,L_z,x_1,x_1'$ determine the
angular phase state modulo this action.  Away from the poles this
follows from $x_1'=-\sin\theta\,p_\theta$; at a pole $D$
determines the nonzero covector's norm, and rotations act transitively
on covectors of that norm.  Therefore there is a unique
$\alpha(z)\in\mathbb S^1$ such that
\begin{equation}
 \Phi^{T_{\mathrm{ang}}(z)}z=R_{\alpha(z)}z.
 \label{geo:angular-return}
\end{equation}
To justify smoothness without treating the listed invariants as
coordinates, let $\mathcal U$ be the angular phase region with
nonzero covectors.  Its free circle action is proper, and the map
\[
 \mathbb S^1\times\mathcal U\longrightarrow
       \mathcal U\times_{\mathcal U/\mathbb S^1}\mathcal U,
 \qquad (\zeta,z)\longmapsto(z,R_\zeta z)
\]
is a diffeomorphism.  The smooth pair
$z\mapsto(z,\Phi^{T_{\mathrm{ang}}(z)}z)$ lies in this fiber
product by the preceding orbit classification.  Its inverse gives
the smooth $\alpha$, including at poles and the equatorial limit.
Invariance of
$T_{\mathrm{ang}}$ and equivariance of the flow under axial rotations
show that $\alpha$ is invariant along the flow.

For every fixed integer $n$, it follows that
\begin{equation}
 \Phi^s z=R_{\alpha(z)^n}
            \Phi^{s-nT_{\mathrm{ang}}(z)}z.
 \label{geo:angular-iterate}
\end{equation}
At each evaluation point choose $n$ so the residual time is bounded,
and keep this integer fixed while differentiating the identity.
On a compact invariant subset the period is bounded above and away
from zero, so $|n|=O(1+|s|)$.  Every fixed derivative of the circle
power map and of $s-nT_{\mathrm{ang}}(z)$ costs only a polynomial
factor in $n$, while bounded-time flow derivatives are uniformly
bounded.  Hence for each fixed integer $j\geq1$ there are constants
$C_j,N_j$ such that
\begin{equation}
 \|D_z^j\Phi^s\|\leq C_j(1+|s|)^{N_j}.
 \label{geo:polynomial-tangent}
\end{equation}
This argument uses the circle-valued return rotation and requires
no global real-valued lift of its angle.

For a positive time change, write
\[
 t=\int_0^s\tau(\Phi^u z)\,du,\qquad 0<c\leq\tau\leq C.
\]
The bounds just proved give polynomial bounds for every fixed
derivative of this integral.  Its $s$-derivative is
$\tau(\Phi^s z)\geq c$; repeated implicit differentiation gives
polynomial bounds for its inverse $s=s(t,z)$, and $|s|$ is
comparable to $|t|$.  Composition with $\Phi^s$ preserves
\eqref{geo:polynomial-tangent}.  The time changes between separated
time and a regular stationary Kerr time have precisely these
properties on the compact trapped set; reverse the separated
orientation on the past sheet when parameterizing by increasing
$t$.

Finally, all formulas hold on an angular neighborhood where
$c_1>0$, allowing $D$ to vary.  On a nearby level of
\eqref{geo:separated-symbol}, the radial critical equations are
unchanged and only the allowed value of $D$ is shifted.  The smooth
embedding $r=r_c(L_z/E)$, $k_r=0$ transfers the estimates to every
tangent jet of the restricted flow on this critical manifold.
\end{proof}

\begin{proposition}[Normally hyperbolic trapping with escape]
\label{geo:normal-hyperbolicity}
For each fixed subextremal Kerr metric, the energy-normalized photon
trap is symplectically transverse of codimension two and has smooth
local stable and unstable surfaces.  Their normal expansion and
contraction rates are bounded away from zero.  All fixed-order
tangent derivatives grow subexponentially, so the normal rates
dominate the tangent dynamics at every fixed order.  These
properties hold on a sufficiently small interval of nearby levels
of the separated fixed-frequency symbol.  A departing radial tail
escapes the radial trapping neighborhood and cannot return along a
homoclinic orbit.
\end{proposition}

\begin{proof}
The functions $\varphi_\pm$ of
Lemma~\ref{geo:normal-functions} give the two invariant surfaces;
their common zero set is the trap, and their nonzero bracket proves
symplectic transversality.  Their differential equations give the
normal exponential rates.  Lemma~\ref{geo:tangent-growth} bounds
every fixed-order tangent derivative polynomially, which proves
the claimed domination.

For nearby levels, use the symbol $q_E$ in
\eqref{geo:separated-symbol} and the spatial field
$X=H_{q_E}/2$.  Continue $z_1$ algebraically as
$z_1=\sqrt\Delta\,k_r/E$.  On every nearby level one has the
exact identity
\begin{equation}
 z_1^2-z_2^2=\frac{q_E}{E^2}+U_c-\frac{D}{E^2}.
 \label{geo:off-null-invariant}
\end{equation}
The quantities $q_E,D,E,L_z$ are constant for $X$, and
$Xr=E\sqrt\Delta\,z_1$.  Therefore
\[
 Xz_2=(\partial_rz_2)E\sqrt\Delta\,z_1,\qquad
 Xz_1=(\partial_rz_2)E\sqrt\Delta\,z_2.
\]
For the full stationary field, differentiation at fixed
Boyer--Lindquist spatial covector gives
$X^T=-\partial_Eq_E/2$.  In the regular time
$t=T+r_*-G$, its speed is instead
\[
 X^t=X^T+(r_*'-G')Xr.
\]
On the trap $Xr=0$, so $E/X^t>0$ there and in a sufficiently
small neighborhood on both sheets.  Dividing by $X^t$ gives the
same positive $\lambda$ for
both equations, with increasing-time orientation.  Thus these
are invariant hypersurfaces in a full fixed-frequency neighborhood,
not merely in its null level.  Their bracket remains nonzero by
continuity.  The positive angular margin also persists, as proved
above.

For a trapped impact parameter, Lemma~\ref{geo:radial} gives
$F(r_+)>0$.  Along a departing tail $L_z/E$ and $D/E^2$ remain
fixed, and $F^2/\Delta$ increases strictly as $r$ moves away from
$r_c$ on either side.  Thus the departure is genuine escape, with
no return to the double root.
\end{proof}

This is the trapping geometry used by the local normally
hyperbolic microlocal estimate below
\cite{Dyatlov2016,Hintz2017}.  The argument requires neither a
pinching-dependent resonance expansion nor a stationary Killing
field that is timelike over the whole trapping base.  The local
stable and unstable tails are the real surfaces just constructed;
outside their neighborhood the radial alternatives permit
propagation to microlocal absorption or to the ordinary ends.
All of these trapping and horizon-source descriptions are invariant
under stationary chart changes.  They also persist, with uniform
positive margins after shrinking the parameter neighborhood, for
small changes of the exact Kerr parameters.

\subsection{Ellipticity after spatial projection}

The remaining issue is specific to arbitrary spin.  A trapping base
point can lie in the ergoregion, so one cannot require the stationary
operator to be elliptic on an entire base neighborhood.  The spatial
projection of a trapped covector nevertheless lies in an elliptic
cone for the slicing of Lemma~\ref{geo:slicing}.

\begin{proposition}[An elliptic cone containing the projected trap]
\label{geo:projected-ellipticity}
Choose $G'=1$ throughout the photon region.  The union of the
spatial projections of both trapped null sheets is contained, with
a positive margin, in a spatial conic neighborhood $\mathcal T$
with compact base support separated from the horizon and from
both spatial ends of the estimates.  On this neighborhood,
\begin{equation}
 b^{ij}(y)\eta_i\eta_j\geq c|\eta|^2
 \label{geo:conic-ellipticity}
\end{equation}
for some $c>0$ and any fixed auxiliary cotangent norm on the
compact base support.  Both temporal roots there are nonzero and
have absolute values comparable to $|\eta|$.  One can choose nested
smaller and larger conic neighborhoods with these properties, and
the same choices work after sufficiently small changes of exact
Kerr parameters and stationary charts.
\end{proposition}

\begin{proof}
At a trapped null covector, $\dot r=0$, so $\dot t=\dot T$.
Substituting \eqref{geo:double-root} into
\eqref{geo:ray-equations} gives
\begin{equation}
 \frac{\Sigma_K k^\sharp(dt)}{E}
   =\frac{(r^2+\mathfrak a^2)F}{\Delta}
        +\mathfrak a\left(\frac{L_z}{E}
                               -\mathfrak a\sin^2\theta\right)
   =\Sigma_K+\frac{4Mr^2}{r-M}.
 \label{geo:trapped-time-speed}
\end{equation}
Write $k=-E\,dt+\eta$ in the stationary splitting.  The
zero-temporal-component lift of its spatial projection is
\[
 \eta^\uparrow=k+E\,dt.
\]
Using nullness, \eqref{geo:time-norm-one}, and
\eqref{geo:trapped-time-speed}, we obtain the exact identity
\begin{equation}
 \begin{aligned}
 \frac{\Sigma_K}{E^2}
       b^{-1}(\eta^\uparrow,\eta^\uparrow)
 &=\frac{2\Sigma_K k^\sharp(dt)}{E}
            +\Sigma_K b^{-1}(dt,dt)\\
 &=\Sigma_K+\frac{8Mr^2}{r-M}-2Mr>0.
 \end{aligned}
 \label{geo:projected-symbol}
\end{equation}
The last inequality is strict for every $r>r_+>M$.

The normalized trap is compact by Lemma~\ref{geo:radial}.
Its spatial projection is nonzero, since a nonzero multiple of
$dt$ is timelike and cannot be a null covector.  Consequently
\eqref{geo:projected-symbol} gives a positive lower bound relative
to $|\eta|^2$ on that projection.  Covector reversal supplies the
same bound on the other sheet.  By continuity, a neighborhood in
the spatial cosphere has the asserted inequality.  Its base
support may be confined to a fixed compact radial interval
separated from the horizon and the spatial ends, and nested
neighborhoods may be chosen inside this open elliptic set.

For completeness, the temporal characteristic equation is
\begin{equation}
 b^{tt}E^2-2b^{ti}\eta_i E+b^{ij}\eta_i\eta_j=0.
 \label{geo:temporal-roots}
\end{equation}
Here $b^{tt}<0$, while the constant term is positive on
$\mathcal T$.  The two real roots therefore have opposite signs
and neither vanishes.  Homogeneity and compactness give the
comparability of their absolute values with $|\eta|$.
The radial description places the trap of every nearby exact Kerr
metric in the same compact neighborhood, after shrinking the
parameter range.  Smooth dependence of the charts, the symbol,
and the trapped set preserves all the strict margins.
\end{proof}

There is no assertion of base-space stationary ellipticity in
Proposition~\ref{geo:projected-ellipticity}.  Elsewhere the temporal
root gap is still comparable to $|\eta|$ on each fixed compact
spatial set: this follows from the timelikeness of $dt$ and the
positive definite inverse metric induced on its slices.  Thus
spatial halfwave diagonalization remains available even where one
temporal root vanishes.  Together with strict outflow and the
escape alternatives, this permits the full-order positive energy
sign away from trapping.  Near infinity the entire zero-frequency
spatial symbol is elliptic.

\begin{lemma}[Zero-energy directions at spatial fiber infinity]
\label{geo:zero-energy}
The zero-temporal-energy characteristic set has no photon trapping
away from the horizon radial set.  Strictly outside the horizon,
its spatial Hamilton flow is real principal and nonradial,
including at a covector dual to the stationary Killing vector on
the stationary limit surface.  Consequently the extra
characteristics at spatial fiber infinity in an ergoregion do not
introduce a further stationary trapped set at bounded temporal
frequency.
\end{lemma}

\begin{proof}
For a future null covector with $E=0$ and $r>r_+$, positivity
of $P=-\mathfrak a L_z$ implies $L_z\neq0$ and hence $D>0$.
The radial polynomial and its derivative are
\begin{equation}
 R(r)=\mathfrak a^2L_z^2-\Delta D,\qquad
 R'(r)=-2(r-M)D<0.
 \label{geo:zero-energy-radial}
\end{equation}
Thus there is at most one radial turn, and it is simple.

The spatial fixed-energy Hamilton field is the projection of the
full null Hamilton field.  If its spatial base component vanishes,
the full metric dual must be proportional to $\partial_T$.
Nullness then places the base point on the stationary limit
surface.  At such a point \eqref{geo:zero-energy-radial} still
has a simple root, so the spatial Hamilton field is nonzero.
It cannot be a nonzero pure fiber dilation: axial momentum
$L_z\neq0$ is conserved, whereas a nonzero dilation would change
it.  This proves the nonradial real-principal assertion even at
these apparently stationary covectors.

Between the horizons, the radial component of a null ray is
nonzero because $dr$ is timelike.  At the horizon, all crossings
other than the generators are transverse.  Lemmas~\ref{geo:radial}
and \ref{geo:horizon-source} therefore exhaust the possible
non-escaping behavior.  At bounded temporal frequency, spatial
fiber infinity corresponds to normalized temporal energy zero,
so it contributes exactly the directions just analyzed and no
additional trapping.
\end{proof}

After the standard horizon-source propagation, with the supported
reverse estimate for the transpose problem, these zero-energy
directions use ordinary real-principal propagation.  More
generally, fixed-energy spatial propagation follows the projected
null flow, oriented by increasing regular time or reversed for
the opposite estimate.  The ray alternatives connect the trap,
the horizon radial region, the far end, and inner outflow without
requiring a positive Killing-current norm in the ergoregion.
Stationary affine-in-time coordinate changes preserve this
description.  These are the geometric inputs for the stationary
and compact estimates in the next sections.
